{\section{The First Experimental Verdict: The Freedman and Clauser Setup}

Bell's generalization, often called the CHSH inequality, was the key theoretical tool that enabled the design of the Freedman and Clauser experiment. 
They demonstrated that the measured coincidence rates had to obey the following constraint:
$$ 1 \le \frac{3R(\varphi)}{R_{0}} - \frac{R(3\varphi)}{R_{0}} - \frac{R_{1}+R_{2}}{R_{0}} \le 0 $$
where $R(\phi)$ is the coincidence rate, $R_0$ is the rate with both polarizers removed, and $R_1$ and $R_2$ are the rates with a single polarizer removed.

To simplify the test, they focused on the angles of predicted maximum violation, $\varphi = 22.5^{\circ}$ and $3\varphi = 67.5^{\circ}$. At these points, 
the inequality can be combined into a much simpler expression:
$$ \delta = \left|\frac{R(22.5^{\circ})}{R_{0}} - \frac{R(67.5^{\circ})}{R_{0}}\right| - \frac{1}{4} \le 0 $$
Local realism demanded that $\delta \le 0$. Quantum mechanics predicted a positive value for $\delta$.

\subsection{The Decisive Experimental Result}

The experimental results for the normalized coincidence rates were as follows:
\begin{itemize}
    \item At the relative angle of 22.5°: $R(22.5^{\circ})/R_{0} = 0.400 \pm 0.007$
    \item At the relative angle of 67.5°: $R(67.5^{\circ})/R_{0} = 0.100 \pm 0.003$
\end{itemize}
Substituting these values into the inequality formula, the experiment yielded the final result:
$$ \delta = |0.400 - 0.100| - 0.25 = 0.300 - 0.25 = 0.050 $$
Including the experimental uncertainties, the final reported value was:
$$ \delta = 0.050 \pm 0.008 $$
This result violates Bell's inequality and confirms the predictions of quantum mechanics.
